\section{Errors of Conception}

\begin{quotex}
Anyone who wants to benefit from the Law of Exception must first achieve a victory over himself, over his own interior world, before he will be able to overcome the `World' and --- by doing so --- escape from the General Law.
\flright{\textsc{Boris Mouravieff}}
\end{quotex}


Conception is defined as the ``imagination or self-originated concept''. So an \textbf{error of conception} is a false conception of the world or of oneself.

\textbf{Thomas Hora}, the Hungarian psychiatrist and founder of existential metapsychiatry, describes five errors of conception, which he colorfully names the ``five gates of hell''.

\begin{enumerate}


\item \textbf{Sensualism}. This mode of being-in-the-world where the primary preoccupation is with sensory awareness, pleasure, and pain. It judges reality purely by sensory awareness. It is concerned with the pursuit of pleasure and the avoidance of pain. ``Feeling good'' is its goal. The sensualist likes or dislikes.

This is Man \#1.

\item \textbf{Emotionalism}. This mode of being-in-the-world is preoccupied with emotional experiences and seeks to cognize reality on the basis of feelings. Emotionalists try to user their feelings as sources of information concerning what is and what is not real. Feelings provide misinformation yet seem to be very valid. The emotionalist is convinced that things are really the way they feel. He is a reactor, not a responder. He is often a sentimentalist and is on a quest for excitement.

This is Man \#2.

\item \textbf{Intellectualism}. This mode of being-in-the-world places great importance on being known as knowing. The intellectualist is a living filing cabinet and likes to display its contents. Often, he is dominated by various impossible conspiracy theories. He is often subject to an ``\emph{idee fixe}'' that prevents him from escaping his false worldview. He enjoys getting into debates and this type has become very prominent on the Internet. They use thought as a weapon rather than as a means to reach the truth.

This is Man \#3.

\item \textbf{Personism}. This mode of being-in-the-world is concerned about ``what others are thinking about what we are thinking.'' The personist is centered on man rather than God. His activities and opinions are calculated to get others to think highly of him. Hence, he may make public claims about secret knowledge, remarkable abilities, miracles, and so on.

Mouravieff relates this to esoteric work:

\begin{quotex}
If the seeker starts with a negative approach and a feeling of inferiority and dissatisfaction---approaching the esoteric domain driven by the desire to find in it \emph{personal} and thus \emph{impure} satisfaction for himself, he will not be able to advance very far along this way. If he persists, he will meet with failure. The error of conception made at the start will imperceptibly lead him towards this \emph{mysticism of phenomena}.
\end{quotex}


\item \textbf{Materialism.} This mode of being-in-the-world is dominated by the need for material objects and possessions. The materialist will be overly concerned about brand names and will make sure that others know what brands he uses.
\end{enumerate}


Besides these erroneous conceptions of being-in-the-world, there are errors of conception about one's own self. Mouravieff provides some concrete examples in the Appendix to Volume 1 of \emph{Gnosis}. For example, there is this story:

\begin{quotex}
Let us suppose that a man has been angered by his wife as he rose from sleep. He avoided any further quarrel, but cut himself badly while shaving, and this infuriates him. As a last straw, he finds a flat tyre on his car. Finding nobody to call, he changes the tyre, soiling his hands before going to his office --- and is still boiling with repressed anger and contained violence when he arrives there. Once in his office, he calls his secretary, and is compelled to wait two minutes before she arrives. When she finally enters, it is quite apparent that she was drinking her coffee. Asked for a special file, she answers that it has not yet been completed: that important papers have been given to a co-director. This is enough for him to burst into a fit of uncontrollably violent rage.
\end{quotex}


He justifies his irrational fit as caused by her alleged insouciance and is unaware of all the preceding factors that created the bad mood in the first place. That is why, in esoteric work, we focus on the expression of negative emotions. In that way, we try to become aware of all the thoughts and events that led up to that expression.

Often a person will try to justify his or her behavior. They will always have an answer to the question, ``Why did you do it?'' but will eventually reach the point where there is no longer an answer. That is the esoteric meaning behind the notion that man cannot justify himself. Rather it must come from a transcendent source. Mouravieff explains:

\begin{quotationx}
Sin is only the expression of an error of conception translated into action, that is to say, of an entirely false attitude to the problems or questions which, in general or in particular, rise before us daily. This is a subject where we often confuse cause and effect. Sin is simply the effect of a causal attitude which necessarily leads to deviation and straying with all their consequences.

Thus the basis of sin is \emph{error.} That is why sin can and must be redeemed. The way of redemption is simple, but how difficult it is to put into practice! The method is \emph{repentance.}

To repent is to become conscious of the error which led to the act of sin.
\end{quotationx}


We see that the way to expose the error of conception is to become conscious of it, not by creating an ``opposite thought''; in our esoteric work, we can see how that will dissipate the error.

Thomas Hora likewise rejects any attempt at self-justification. Rather, it is necessary to grasp the meaning behind. He suggests asking oneself two ``intelligent questions'':

\begin{itemize}


\item What is the meaning of what seems to be?

\item What is what really is?
\end{itemize}


By grasping the meaning of a situation, we can understand what is really going on. Meaning is not found by being-in-the-world, but rather by transcending it and relating reality to the Real I. You may recall what \textbf{Julius Evola} wrote in regard to \textbf{Hermann Keyserling}:


\begin{quotex} 
Keyserling contemplates the point in which the I says to himself: ``I understood.'' It is essentially a point of spontaneity, freedom, and interiority: there is no way to compel anyone to understand. It on the other hand has a mystico-illuminative character: it seeks to live the moment of ``meaning'', not of this or that meaning, but rather of meaning in general, of the pure element of the understanding conditioning every understanding, then one will feel that something ineffable shines in it, something that, even if contending with it and being supported on it, absolutely transcends the whole of the means and forms from which it was propitiated. This mystical moment of pure understanding is the moment of the spirit.
\end{quotex}

\textit{References:}

Thomas Hora, \emph{Beyond the Dream}

Boris Mouravieff, \emph{Gnosis}


\flrightit{Posted on 2014-10-24 by Cologero}

\begin{center}* * *\end{center}

\begin{footnotesize}
\begin{sffamily}
\texttt{August on 2014-10-25 at 08:20 said:}

What about the case of one who is aware of preceding factors but decides to have a fit anyway? If some of us have used esoteric techniques to develop awareness but are not in a position to quit the world, why should we be patient with fools? Why go to the effort of using awareness as a circuit breaker that trips irrational fits in an irrational social life?

The only good reason I can think of is avoiding the consequences, but in many cases they are not severe enough so I express my displeasure directly, and would perhaps do so more often in an environment where it were beneficial.

\hfill

\texttt{Cologero on 2014-10-27 at 01:06 said:}

August, perhaps you can go back and read \textit{We Anti-moderns}\footnote{\url{http://www.gornahoor.net/?p=3705}} by Julius Evola where we wrote:

\begin{quote}\footnotesize

Haven't we had enough kvetching and jeremiads? We know things are awry, so what is the point of indulging in criticisms of this injustice or that stupid policy? There is a time for righteous anger and a time for sober reflection. Pace this critique, there must be men who stand above them.

\begin{itemize}


\item They know the ancient certitudes.

\item They are certain about their principles.

\item They seek to recover their traditions.

\item Their spirits are unified in solidarity each other and continuity with their ancestors.

\item They are aware of the dark and chaotic forces within themselves and learn to master them.
\end{itemize}

\end{quote}


Nevertheless, there is certainly just cause for \textit{righteous anger}\footnote{\url{http://www.gornahoor.net/?p=5471}}.


But lashing out that has no purpose is not indicative of self-mastery or holiness:

\begin{quote}\footnotesize

We must clean out our hearts, ridding them of darkness and bitterness; we must make them clean and sparkling places for God to live. We must be as thorough with ourselves as we are with our homes. Are there any dark corners in our hearts we have avoided for so long? Are we simply ``sweeping all the dirt under the rug?'' God sees all and knows all. He knows what is behind every wall of our hearts, what is swept into every corner, and what is hidden under every rug. Let us truly clean out our hearts; let us rid ourselves of the grudges, pain, and \textbf{anger} that clutter our ability to love freely. Let us empty out every nook and cranny, so that His divine light can shine throughout!'' 
\flright{\textsc{Swami Chidanand Saraswati}}
\end{quote}


\hfill

\texttt{August on 2014-10-27 at 03:49 said:}

Thanks for those reminders Cologero, you're quite right, and it's clear sight of the inner work that should give us pause and a place to retreat when lashed by insults from without. The various obstacles that deliberately try to draw us out of, or stall our return to, this refuge should be seen for what they are.

Some of us know of the custom that in your parts is called Vendetta, and all the more so if we actively worship our forebears. I think it wasn't only a habit of the subordinate castes. But so it is -- in these times let us focus on self-mastery and holiness, those of us who are destined to suffer for it must do so one way or another, for a time, paying a price, sometimes uncertain.

Especially let the young ones be mindful of this as the glow of youth begins to recede and one is faced with the harder reality of his choices, commitments, circumstances.

\hfill

\texttt{August on 2014-10-27 at 04:34 said:}

Pardon me, semantic clarification: to avoid misreadings, in my comment below the sentence starting ``But so it is... '' should be taken as ``But anyway, in these times let us focus on... ''

The way it currently appears, it could be read as reversing the statement in the preceding sentence, which is not the intended meaning.


\hfill

\end{sffamily}
\end{footnotesize}

